\documentclass[10pt,a4paper]{amsart}
\usepackage[margin=19mm]{geometry}
\usepackage{amsmath,amssymb,mathtools}
\usepackage[scr=rsfs,cal=euler]{mathalfa}
\usepackage{xfrac}
\usepackage[unicode,hidelinks,bookmarks=false]{hyperref}
\newcommand{\ie}{i.e.\ }
\newcommand{\cf}{cf.\ }
\newcommand{\resp}{respectively\ }
\newcommand{\Subsection}{\subsection}
\providecommand{\nobreakdash}{\nobreak\hbox{-}}
\setlength{\emergencystretch}{2em}
\multlinegap0pt
\usepackage{amssymb}
\newtheorem{theorem}{Theorem}
\newcommand{\CC}{\mathbb C}
\newcommand{\K}{\mathcal K}
\title{On the connection between coordinate and diagonal arrangement complements}
\author{Vsevolod Tril}
\date{Source: arXiv:2409.18001v1 (CC BY 4.0)}
\begin{document}
\maketitle

\noindent\emph{Source and licence.} On the connection between coordinate and diagonal arrangement complements. Version arXiv:2409.18001v1 (CC BY 4.0), distributed by arXiv under the Creative Commons Attribution 4.0 International licence, http://creativecommons.org/licenses/by/4.0/, as recorded in the arXiv licence metadata for this exact version and shown on the abstract page https://arxiv.org/abs/2409.18001v1. Full licence notice: \texttt{licenses/arxiv-2409.18001-CC-BY-4.0.txt}.\par\smallskip

\noindent\textit{Editorial scope.} Counted unit: the article's theorem res1 (source0.tex lines 406--412). The article's complete original proof of it is delivered with it (source0.tex lines 413--448, one \texttt{proof} environment of 36 lines). No mathematics is written, repaired or re-derived: the statement and the proof are the article's own lines, in the article's own notation, and no retained line is substituted or re-flowed.  No further source-local context is needed: every cross-reference of every retained line resolves inside the two delivered spans.  Nothing was omitted from any retained span, so the package carries no omission record.  No retained line contains a citation, so the wrapper carries no bibliography and no bibliography entry.  No retained line contains a figure environment, an image-inclusion command or drawing code, so no image is delivered and the package declares no asset dependency.  Editorial formatting only, in addition: an AMS \texttt{amsart} preamble that restates the article's own declarations and macros used by the retained lines, namely the environments \texttt{theorem} on separate counters, together with the standard packages those lines need.  The retained environments print in this document as Theorem 1 in order of appearance, whereas the article prints the same items with the numbering of its own full document.  The complete native source of the article as distributed in the arXiv e-print for this exact version is bundled unchanged as \texttt{sources/165866/source0.tex}.
\begin{theorem}\label{res1}
There is a homotopy equivalence
%	Let $\K$ be a simplicial complex on $[m]$.
%	Then there is a simplicial complex $\mathcal{L}$ on the vertex set $[m+1]$ such that
	$$U(\K) \simeq D(\mathcal{L}).$$
The same is true for real arrangement complements.
\end{theorem}

\begin{proof}
	The map
 %Define the following map:
	$$F: \CC^{m+1}\times [0,1] \to \CC^{m+1},$$
	$$F(z_1, \dots, z_{m+1}, t) = (z_1, \dots, z_{m+1}) - t
	\frac{z_1 + \dots + z_{m+1}}{m+1} (1, \dots, 1)$$
	defines a homotopy between the identity map $id: \CC^{m+1} \to \CC^{m+1}$ and the orthogonal projection
	$P: \CC^{m+1} \to \CC^{m+1}$ onto the hyperplane $\pi = \{z_1 + \dots + z_m = 0\}$.
	
Observe that $F(D(\mathcal{L}), t) \subset D(\mathcal{L})$ for every $t \in [0, 1]$. 
	Indeed, the equality $z_{i_1} = \dots = z_{i_k}$ holds if and only if
	$$z_{i_1} - t \frac{z_1 + \dots + z_{m+1}}{m+1} = \dots = z_{i_k} - t \frac{z_1 + \dots + z_{m+1}}{m+1}.$$
	Hence $z \notin D(\mathcal{L})$ implies $z \notin F(D(\mathcal{L}), t)$, that is $F(D(\mathcal{L}), t) \subset D(\mathcal{L})$.
	We also have that $F(z, t) = z$ for any $z \in \pi$, $t \in [0, 1]$,
	so we obtain $F(D(\mathcal{L}), 1) = D(\mathcal{L}) \cap \pi$,
	and therefore $D(\mathcal{L}) \simeq D(\mathcal{L}) \cap \pi$.
	
	Consider the linear isomorphism $A: \CC^{m+1} \to \CC^{m+1}$ defined by the formula
	$$y_1 = z_1 - z_{m+1},\quad \dots, \quad y_m = z_m - z_{m+1}, \quad y_{m+1} = z_1 + \dots + z_m + z_{m+1}.$$
 %$$\begin{array}{lcl}
	%	y_i = z_i - z_{m+1}, \text{ if } i = 1, \dots, m, \\
	%	y_{m+1} = z_1 + \dots z_m + z_{m+1}.
	%\end{array}$$
	We claim that $A(D(\mathcal{L}) \cap \pi) = U(\K)$, where $U(\K) \subset \{y_{m+1} = 0\} \subset \CC^{m+1}$. 
	Indeed, %the definition of $\K '$ implies that 
    the non-faces of $\mathcal{L}$ are of the form $\{i_1, \dots, i_k, m+1\}$,
	where $\{i_1, \dots, i_k\} \notin \K$. Therefore,
\begin{multline*}
  A(D(\mathcal{L}) \cap \pi) = A\Bigl( { \CC^{m+1} \backslash\!\!\!\!\bigcup_{\{i_1, \dots, i_k\} \notin \K}
		\{z_{i_1} = \dots = z_{i_k} = z_{m+1} \} }\Bigr)  \cap \,A \bigl({ \{z_1 + \dots + z_{m+1} = 0\} }\bigr) \\
		= \Bigl({\CC^{m+1} \backslash \bigcup_{\{i_1, \dots, i_k\} \notin \K}
		\{y_{i_1} = \dots = y_{i_k} = 0 \} }\Bigr) \cap \{y_{m+1} = 0 \} = U(\K),
\end{multline*}  
where the last equality holds since $U(\K)$ is considered as a subspace in the hyperplane $\{y_{m+1} = 0\}$.
It follows that $A\colon \CC^{m+1} \to \CC^{m+1}$ induces a homeomorphism between $D(\mathcal{L}) \cap \pi$ and $U(\K)$, so we obtain a homotopy equivalence $D(\mathcal{L}) \simeq D(\mathcal{L})\cap \pi \cong U(\K)$.
\end{proof}

\end{document}
